\unitchapter{Paper 1 · Unit 7}{Data Interpretation}{p1-07}

\textbf{Expected questions: 5 (one data set) · Target: 5/5 --- pure calculation marks}

\begin{sylbox}

\begin{itemize}
\tightlist
\item[$\square$]
  Sources, acquisition and classification of data
\item[$\square$]
  Quantitative and qualitative data
\item[$\square$]
  Graphical representation (bar chart, histograms, pie chart, table chart, line chart) and mapping of data
\item[$\square$]
  Data interpretation
\item[$\square$]
  Data and governance
\end{itemize}

\end{sylbox}

\needdiagram{figures/p1-07-00}{55pt}

\hypertarget{p1-07--1-data-basics}{%
\section{1. Data Basics}\label{p1-07--1-data-basics}}

\diagramhere{figures/p1-07-00}

\begin{itemize}
\tightlist
\item
  Classification bases: \textbf{Geographical (spatial), Chronological (temporal), Qualitative, Quantitative}.
\item
  Data mapping: choropleth maps (shades by region), dot maps, isopleths, cartograms.
\end{itemize}

\needspace{135pt}

\hypertarget{p1-07--2-graph-types--when-to-use}{%
\section{2. Graph Types --- When to Use}\label{p1-07--2-graph-types--when-to-use}}

\begingroup\small

\begin{longtable}[]{@{}
  >{\raggedright\arraybackslash}p{(\columnwidth - 4\tabcolsep) * \real{0.1466}}
  >{\raggedright\arraybackslash}p{(\columnwidth - 4\tabcolsep) * \real{0.3456}}
  >{\raggedright\arraybackslash}p{(\columnwidth - 4\tabcolsep) * \real{0.2359}}@{}}
\toprule\noalign{}
\begin{minipage}[b]{\linewidth}\raggedright
\textbf{Graph}
\end{minipage} & \begin{minipage}[b]{\linewidth}\raggedright
\textbf{Best for}
\end{minipage} & \begin{minipage}[b]{\linewidth}\raggedright
\textbf{Note}
\end{minipage} \\
\midrule\noalign{}
\endhead
\bottomrule\noalign{}
\endlastfoot
Bar chart & Comparing categories & Gaps between bars \\
Histogram & Frequency of continuous classes & \textbf{No gaps}; area ∝ frequency \\
Frequency polygon & Joining midpoints of histogram & \\
Ogive & Cumulative frequency (less than / more than) & Intersection gives \textbf{median} \\
Pie chart & Part of a whole & Angle = (value/\allowbreak{}total) × 360° \\
Line chart & Trend over time & \\
Scatter plot & Correlation between two variables & \\
Table & Exact values & \\
\end{longtable}

\endgroup

\begin{asciiblock}{57pt}
\begin{Verbatim}[fontsize=\fontsize{8.8}{10.7}\selectfont]
Pie chart conversions
 1% = 3.6°      25% = 90°      50% = 180°
 angle θ  ⇒  percentage = θ / 3.6
\end{Verbatim}
\end{asciiblock}

\needspace{135pt}

\hypertarget{p1-07--3-measures-of-central-tendency--dispersion}{%
\section{3. Measures of Central Tendency \& Dispersion}\label{p1-07--3-measures-of-central-tendency--dispersion}}

\begingroup\small

\begin{longtable}[]{@{}
  >{\raggedright\arraybackslash}p{(\columnwidth - 4\tabcolsep) * \real{0.1869}}
  >{\raggedright\arraybackslash}p{(\columnwidth - 4\tabcolsep) * \real{0.2097}}
  >{\raggedright\arraybackslash}p{(\columnwidth - 4\tabcolsep) * \real{0.2117}}@{}}
\toprule\noalign{}
\begin{minipage}[b]{\linewidth}\raggedright
\textbf{Measure}
\end{minipage} & \begin{minipage}[b]{\linewidth}\raggedright
\textbf{Formula}
\end{minipage} & \begin{minipage}[b]{\linewidth}\raggedright
\textbf{Note}
\end{minipage} \\
\midrule\noalign{}
\endhead
\bottomrule\noalign{}
\endlastfoot
Mean & Σx / n & Affected by extreme values \\
Median & Middle value & Best for skewed data \\
Mode & Most frequent & Only one for nominal data \\
Empirical relation & Mode = 3 Median − 2 Mean & Moderately skewed \\
Range & Max − Min & \\
Variance & Σ(x − x̄)² / n & \\
SD & √Variance & \\
Coefficient of variation & SD / Mean × 100 & Compare consistency \\
\end{longtable}

\endgroup

\needspace{144pt}

\hypertarget{p1-07--4-speed-calculation-tricks}{%
\section{4. Speed-Calculation Tricks}\label{p1-07--4-speed-calculation-tricks}}

\begin{asciiblock}{89pt}
\begin{Verbatim}[fontsize=\fontsize{8.8}{10.7}\selectfont]
Percentage change    = (New − Old)/Old × 100
Percentage points    = simple difference of two percentages
Ratio comparisons    → compare cross-products: a/b > c/d  ⇔ ad > bc
CAGR                 = (End/Start)^(1/n) − 1
Fraction-percent     1/2=50 1/3=33.33 1/4=25 1/5=20 1/6=16.67 1/7=14.28
                     1/8=12.5 1/9=11.11 1/11=9.09 1/12=8.33 1/16=6.25
\end{Verbatim}
\end{asciiblock}

Approximate first, then pick nearest option; options are usually far apart.

\needspace{170pt}

\hypertarget{p1-07--5-worked-data-set-typical-net-format}{%
\section{5. Worked Data Set (typical NET format)}\label{p1-07--5-worked-data-set-typical-net-format}}

Production of cars (thousands) by 4 companies:

\begingroup\small

\begin{longtable}[]{@{}
  >{\raggedright\arraybackslash}p{(\columnwidth - 10\tabcolsep) * \real{0.0589}}
  >{\raggedright\arraybackslash}p{(\columnwidth - 10\tabcolsep) * \real{0.0372}}
  >{\raggedright\arraybackslash}p{(\columnwidth - 10\tabcolsep) * \real{0.0372}}
  >{\raggedright\arraybackslash}p{(\columnwidth - 10\tabcolsep) * \real{0.0372}}
  >{\raggedright\arraybackslash}p{(\columnwidth - 10\tabcolsep) * \real{0.0372}}
  >{\raggedright\arraybackslash}p{(\columnwidth - 10\tabcolsep) * \real{0.0666}}@{}}
\toprule\noalign{}
\begin{minipage}[b]{\linewidth}\raggedright
\textbf{Year}
\end{minipage} & \begin{minipage}[b]{\linewidth}\raggedright
\textbf{P}
\end{minipage} & \begin{minipage}[b]{\linewidth}\raggedright
\textbf{Q}
\end{minipage} & \begin{minipage}[b]{\linewidth}\raggedright
\textbf{R}
\end{minipage} & \begin{minipage}[b]{\linewidth}\raggedright
\textbf{S}
\end{minipage} & \begin{minipage}[b]{\linewidth}\raggedright
\textbf{Total}
\end{minipage} \\
\midrule\noalign{}
\endhead
\bottomrule\noalign{}
\endlastfoot
2019 & 40 & 30 & 25 & 15 & 110 \\
2020 & 45 & 35 & 20 & 20 & 120 \\
2021 & 50 & 30 & 30 & 25 & 135 \\
2022 & 60 & 40 & 35 & 25 & 160 \\
\end{longtable}

\endgroup

\begin{asciiblock}{78pt}
\begin{Verbatim}[fontsize=\fontsize{8.8}{10.7}\selectfont]
Bar view of P across years
2019 ████████████████████ 40
2020 ██████████████████████▌ 45
2021 █████████████████████████ 50
2022 ██████████████████████████████ 60
\end{Verbatim}
\end{asciiblock}

\begin{enumerate}
\def\labelenumi{\arabic{enumi}.}
\tightlist
\item
  \% increase in total from 2019 to 2022 = 50/110 × 100 = \textbf{45.45\%}
\item
  Average production of Q = 135/4 = \textbf{33.75}
\item
  In which year did R have the highest share of total? 2019: 22.7\%, 2020: 16.7\%, 2021: 22.2\%, 2022: 21.9\% → \textbf{2019}
\item
  Ratio of P(2022) to S(2019+2020) = 60 : 35 = \textbf{12 : 7}
\item
  Pie-chart angle for S in 2021 = 25/135 × 360 = \textbf{66.67°}
\end{enumerate}

\needspace{100pt}

\hypertarget{p1-07--6-data-and-governance}{%
\section{6. Data and Governance}\label{p1-07--6-data-and-governance}}

\begin{itemize}
\tightlist
\item
  \textbf{Open Government Data (OGD) Platform India} --- data.gov.in (NDSAP 2012 --- National Data Sharing \& Accessibility Policy).
\item
  \textbf{Digital Personal Data Protection Act, 2023 (DPDP)} --- rights of Data Principal, duties of Data Fiduciary, Data Protection Board.
\item
  \textbf{National Data Governance Framework Policy (2022)}, India Data Management Office.
\item
  \textbf{Aadhaar (UIDAI, 2009; Aadhaar Act 2016)}, \textbf{e-governance}: Digital India (2015), DigiLocker, UMANG, GSTN, PFMS.
\item
  Data governance principles: quality, privacy, security, accountability, transparency, interoperability.
\item
  Census of India --- every 10 years (Registrar General \& Census Commissioner, MHA); NSO (merged NSSO + CSO, 2019) under MoSPI.
\end{itemize}

\needdiagram{figures/p1-07-01}{130pt}

\hypertarget{p1-07--7-deeper-dive--pie-chart--line-graph-sets-common-pitfalls}{%
\section{7. Deeper Dive --- Pie Chart \& Line Graph Sets, Common Pitfalls}\label{p1-07--7-deeper-dive--pie-chart--line-graph-sets-common-pitfalls}}

\needspace{135pt}

\hypertarget{p1-07--71-pie-chart-set-worked}{%
\subsection{7.1 Pie-Chart Set (worked)}\label{p1-07--71-pie-chart-set-worked}}

A university\textquotesingle s annual expenditure of \textbf{₹80 crore} is distributed as follows:

\diagramhere{figures/p1-07-01}

\begingroup\small

\begin{longtable}[]{@{}
  >{\raggedright\arraybackslash}p{(\columnwidth - 6\tabcolsep) * \real{0.1315}}
  >{\raggedright\arraybackslash}p{(\columnwidth - 6\tabcolsep) * \real{0.0350}}
  >{\raggedright\arraybackslash}p{(\columnwidth - 6\tabcolsep) * \real{0.1662}}
  >{\raggedright\arraybackslash}p{(\columnwidth - 6\tabcolsep) * \real{0.1281}}@{}}
\toprule\noalign{}
\begin{minipage}[b]{\linewidth}\raggedright
\textbf{Head}
\end{minipage} & \begin{minipage}[b]{\linewidth}\raggedright
\textbf{\%}
\end{minipage} & \begin{minipage}[b]{\linewidth}\raggedright
\textbf{Amount (₹ crore)}
\end{minipage} & \begin{minipage}[b]{\linewidth}\raggedright
\textbf{Central angle}
\end{minipage} \\
\midrule\noalign{}
\endhead
\bottomrule\noalign{}
\endlastfoot
Salaries & 45 & 36 & 162° \\
Research & 20 & 16 & 72° \\
Infrastructure & 15 & 12 & 54° \\
Library \& ICT & 10 & 8 & 36° \\
Scholarships & 10 & 8 & 36° \\
\end{longtable}

\endgroup

\begin{enumerate}
\def\labelenumi{\arabic{enumi}.}
\tightlist
\item
  Research exceeds Library \& ICT by: 16 − 8 = \textbf{₹8 crore} (= 100\% more).
\item
  If salaries rise by 10\% and total stays the same, salaries\textquotesingle{} new share = 39.6/80 = \textbf{49.5\%}.
\item
  Ratio of infrastructure to scholarships = \textbf{3 : 2}.
\item
  Angle for research + scholarships = \textbf{108°}.
\item
  If research grants come only from a ₹12 crore grant plus internal funds, internal funds = 16 − 12 = \textbf{₹4 crore}.
\end{enumerate}

\needspace{211pt}

\hypertarget{p1-07--72-two-series-set-worked}{%
\subsection{7.2 Two-Series Set (worked)}\label{p1-07--72-two-series-set-worked}}

Number of PhD scholars admitted (two departments):

\begin{asciiblock}{121pt}
\begin{Verbatim}[fontsize=\fontsize{8.8}{10.7}\selectfont]
 Scholars admitted (each ▇ = 5 scholars)
 2021  CS    ▇▇▇▇▇▇            30
       Maths ▇▇▇▇              20
 2022  CS    ▇▇▇▇▇▇▇▇          40
       Maths ▇▇▇▇▇▇            30
 2023  CS    ▇▇▇▇▇▇▇▇▇▇        50
       Maths ▇▇▇▇▇▇▇           35
 2024  CS    ▇▇▇▇▇▇▇▇▇▇▇▇      60
       Maths ▇▇▇▇▇▇▇▇▇         45
\end{Verbatim}
\end{asciiblock}

\begingroup\small

\begin{longtable}[]{@{}
  >{\raggedright\arraybackslash}p{(\columnwidth - 6\tabcolsep) * \real{0.0553}}
  >{\raggedright\arraybackslash}p{(\columnwidth - 6\tabcolsep) * \real{0.0452}}
  >{\raggedright\arraybackslash}p{(\columnwidth - 6\tabcolsep) * \real{0.0643}}
  >{\raggedright\arraybackslash}p{(\columnwidth - 6\tabcolsep) * \real{0.0626}}@{}}
\toprule\noalign{}
\begin{minipage}[b]{\linewidth}\raggedright
\textbf{Year}
\end{minipage} & \begin{minipage}[b]{\linewidth}\raggedright
\textbf{CS}
\end{minipage} & \begin{minipage}[b]{\linewidth}\raggedright
\textbf{Maths}
\end{minipage} & \begin{minipage}[b]{\linewidth}\raggedright
\textbf{Total}
\end{minipage} \\
\midrule\noalign{}
\endhead
\bottomrule\noalign{}
\endlastfoot
2021 & 30 & 20 & 50 \\
2022 & 40 & 30 & 70 \\
2023 & 50 & 35 & 85 \\
2024 & 60 & 45 & 105 \\
\end{longtable}

\endgroup

\begin{enumerate}
\def\labelenumi{\arabic{enumi}.}
\tightlist
\item
  Year with the highest \% growth for Maths: 2022 (50\%), 2023 (16.7\%), 2024 (28.6\%) → \textbf{2022}.
\item
  Average CS admissions = 180/4 = \textbf{45}.
\item
  Total growth 2021→2024 = (105 − 50)/50 = \textbf{110\%}.
\item
  Ratio CS : Maths over four years = 180 : 130 = \textbf{18 : 13}.
\item
  CAGR of CS (2021→2024, 3 years) = (60/30)\^{}(1/3) − 1 ≈ \textbf{26\%}.
\end{enumerate}

\needspace{100pt}

\hypertarget{p1-07--73-pitfalls-examiners-exploit}{%
\subsection{7.3 Pitfalls Examiners Exploit}\label{p1-07--73-pitfalls-examiners-exploit}}

\begin{itemize}
\tightlist
\item
  \textbf{Percentage vs percentage points}: 20\% → 25\% is a rise of 5 percentage points but 25\%.
\item
  \textbf{Base year confusion}: "\% increase from 2022 to 2023" uses 2022 as the base.
\item
  \textbf{Average of percentages} ≠ overall percentage (weights differ).
\item
  \textbf{Units}: lakhs vs crores (1 crore = 100 lakh); thousands in the table header.
\item
  Read the \textbf{question first}, then compute only what is asked.
\end{itemize}

\needspace{100pt}

\hypertarget{p1-07--74-data-mapping--governance-extras}{%
\subsection{7.4 Data Mapping \& Governance Extras}\label{p1-07--74-data-mapping--governance-extras}}

\begin{itemize}
\tightlist
\item
  \textbf{GIS} layers data on maps; \textbf{choropleth} uses shading per region; \textbf{isopleth} lines join equal values (isotherms, isobars).
\item
  \textbf{Data lifecycle}: collection → storage → processing → analysis → sharing → archival/deletion.
\item
  \textbf{Metadata} = data about data. \textbf{Open data} principles: accessible, machine-readable, licence-free.
\end{itemize}

\needspace{135pt}

\hypertarget{p1-07--previous-year-questions-pyq-pattern}{%
\section{Previous Year Questions (PYQ pattern)}\label{p1-07--previous-year-questions-pyq-pattern}}

DI questions always come as a set of 5 on one table/graph. Typical stems:

\begin{enumerate}
\def\labelenumi{\arabic{enumi}.}
\tightlist
\item
  "What is the percentage increase in X from year A to year B?"
\item
  "In which year was the ratio of X to Y the maximum?"
\item
  "What is the average of X over the given years?"
\item
  "The central angle corresponding to X in the pie chart is\ldots"
\item
  "What is the difference between total X and total Y?"
\item
  "For how many years was X above its average?"
\end{enumerate}

Concept questions:

\begin{enumerate}
\def\labelenumi{\arabic{enumi}.}
\setcounter{enumi}{6}
\tightlist
\item
  A histogram is used for: \textbf{continuous frequency distribution}
\item
  The median can be obtained graphically from: \textbf{Ogives (intersection of less-than and more-than)}
\item
  Data collected by the researcher for the first time: \textbf{Primary data}
\item
  Census data used by a researcher is: \textbf{Secondary data}
\item
  Pie chart: if a sector is 72°, it represents \textbf{20\%}
\item
  Classification of data by time is called: \textbf{Chronological}
\item
  Which measure is best for ordinal data? \textbf{Median}
\item
  Which portal provides open government data in India? \textbf{data.gov.in}
\item
  Coefficient of variation is used to compare: \textbf{consistency/variability of two series}
\end{enumerate}

\textbf{More practice questions}

\begin{enumerate}
\def\labelenumi{\arabic{enumi}.}
\setcounter{enumi}{15}
\tightlist
\item
  A pie-chart sector of 54° represents what \%? \textbf{15\%}
\item
  If total = ₹500 crore and a sector is 18\%, amount = \textbf{₹90 crore}
\item
  Rise from 40\% to 50\% is how many percentage points? \textbf{10}, and what \% increase? \textbf{25\%}
\item
  Average of 20, 30, 40 with frequencies 2, 3, 5 = (40 + 90 + 200)/10 = \textbf{33}
\item
  Lines joining places of equal rainfall on a map: \textbf{isohyets} (a type of isopleth)
\item
  A map using colour shades by state for literacy rate: \textbf{choropleth map}
\end{enumerate}

\begin{revbox}

\begin{itemize}
\tightlist
\item
  1\% = 3.6° · Mode = 3Median − 2Mean · CV = SD/Mean × 100
\item
  Histogram no gaps; bar chart gaps · Ogive → median
\item
  Learn fraction ↔ percent table by heart
\end{itemize}

\end{revbox}
